% ------------------------------------------------------------
% Page 78
% ------------------------------------------------------------

\begin{center}
\begin{minipage}{0.55\textwidth}

On venait l’admirer, ce noble et fier géant,\\
Les touristes lassés, s’asseyaient sur son banc\\
Placé près du vieux pont, pour contempler à l’aise\\
Le site séduisant du vallon de la Brèze.

\vspace{0.35cm}

Nous l’aimions, il était le vaste souvenir\\
De notre enfance heureuse, et n’aurait dû mourir\\
Que longtemps après nous, pour parler de nos vies\\
Aux fils de nos enfants, jouant dans sa prairie

\vspace{0.35cm}

Mais un passant tardif, du sort, vil ouvrier\\
A débroussé sa pipe au pied du poivrier\\
L’étincelle a gagné le cœur du sycomore\\
Et le feu l’a rongé, subtil, jusqu’à l’aurore.

\vspace{0.35cm}

Il était beau ! Qu’importe, il lui fallait périr\\
Il était grand, hélas quand la voix implacable\\
Du trépas nous condamne, il nous faut obéir\\
Tous, tant grands que petits, à l’ordre inexorable.

\vspace{0.35cm}

Mais au printemps prochain quelques rejets\\
Renaîtront vigoureux de ses cendres immenses,\\
Et le destin vengeur de l’acte du pervers\\
Saura lui ménager une longue existence.

\vspace{0.35cm}

La même sève ardente, en ses mille canaux\\
Coulera sans arrêt pendant plus de cent lustres\\
Et que le vieux poivrier vivra dans les rameaux\\
D’un puissant sycomore, à son tour très illustre

\vspace{0.35cm}

Ainsi, bientôt, pour nous la vie aura pris fin\\
Mais la sève, abandonnant les restes\\
De son corps périssable à son heureux destin\\
Fera revivre un corps glorieux et céleste.

\vspace{0.55cm}

\textit{Meyrueis le 30 Octobre 1928}\\
\textit{Signature presque illisible : un nommé REY ? ? ?}

\end{minipage}
\end{center}

\vfill

\begin{flushright}
\textit{Poème trouvé dans les archives}\\
\textit{de Marceline Causse}\\
\textit{à Pradines}\\[0.15cm]
Octobre 2002
\end{flushright}
